\section{Введение в непрерывные игры}

В отличие от конечных игр, где у каждого игрока имеется конечное число стратегий и гарантированно существует равновесие Нэша в смешанных стратегиях, в непрерывных играх множества стратегий представляют собой континуум. Типичными примерами таких игр являются модели выбора объема производства (модель Курно), выбора местоположения (модель Хотеллинга) или уровня усилий в турнирах. Функция полезности в таких играх зачастую является дифференцируемой по стратегиям игрока.

\begin{definition}[Непрерывная игра]
Игра $G = \langle I, S, u \rangle$ называется непрерывной игрой, если для всех игроков $i \in I$ множество стратегий $S_i$ является выпуклым подмножеством конечномерного евклидова пространства $\mathbb{R}^{d_i}$, а функция выигрышей $u_i(s)$ является непрерывной по профилю стратегий $s$.
\end{definition}

Для установления факта существования равновесия в чистых стратегиях в непрерывных играх ключевую роль играет свойство квазивогнутости функции полезности.

\begin{definition}[Квазивогнутая функция]
Пусть $X$ — выпуклое подмножество конечномерного евклидова пространства. Функция $u : X \to \mathbb{R}$ является квазивогнутой, если для любого значения $\bar{u} \in \mathbb{R}$ верхнее лебегуево множество $\{x \mid u(x) \ge \bar{u}\}$ является выпуклым.
\end{definition}

\begin{theorem}[Существование равновесия Нэша в чистых стратегиях]\label{thm:pure_nash_continuous}
Рассмотрим непрерывную игру $G = \langle I, S, u \rangle$, в которой множества стратегий $S_i$ являются выпуклыми компактами. Предположим, что для каждого игрока $i \in I$ функция выигрышей $u_i(s_i, s_{-i})$ является квазивогнутой по $s_i$ для всех $s_{-i} \in S_{-i}$. Тогда в игре $G$ существует равновесие Нэша в чистых стратегиях.
\end{theorem}

Теорема Нэша для конечных игр является частным случаем данной теоремы, поскольку смешанное расширение конечной игры представляет собой непрерывную игру на симплексе (выпуклом компакте) с линейными (а значит, и квазивогнутыми) функциями ожидаемой полезности. Доказательство Теоремы \ref{thm:pure_nash_continuous} опирается на теорему о неподвижной точке. Следует отметить, что равновесий в непрерывных играх может быть несколько.

Нарушение условия квазивогнутости может приводить к отсутствию равновесия Нэша в чистых стратегиях.

\begin{example}[Отсутствие равновесия при нарушении квазивогнутости]
Пусть множества стратегий двух игроков заданы как $S_1 = S_2 = [0, 1]$. Функции полезности имеют вид:
\[
u_1(s_1, s_2) = |s_1 - s_2|, \quad u_2(s_1, s_2) = -s_2^2 + \left(\frac{1}{2} + s_1\right) s_2
\]
Функция $u_1$ не является квазивогнутой по $s_1$. Вычисление наилучших ответов (best responses) дает:
\[
\BR_1(s_2) = 
\begin{cases} 
0, & s_2 > \frac{1}{2} \\
\{0, 1\}, & s_2 = \frac{1}{2} \\
1, & s_2 < \frac{1}{2}
\end{cases}
\]
\[
\BR_2(s_1) = 0.25 + \frac{s_1}{2}
\]
Анализ графиков отображений наилучших ответов показывает, что они не пересекаются, следовательно, равновесие Нэша в чистых стратегиях в данной игре отсутствует.
\end{example}

Если равновесие в чистых стратегиях не существует, рассматривается возможность существования равновесия в смешанных стратегиях.

\begin{theorem}[Существование смешанного равновесия Нэша]
Рассмотрим игру $G = \langle I, S, u \rangle$, в которой множества стратегий $S_i$ являются выпуклыми и компактными подмножествами конечномерных евклидовых пространств $\mathbb{R}^{d_i}$. Предположим, что для каждого игрока $i$ функция выигрышей $u_i(s_i, s_{-i})$ является непрерывной по $s = (s_i, s_{-i})$. Тогда в игре $G$ существует равновесие Нэша в смешанных стратегиях.
\end{theorem}


\section{Приложения непрерывных игр}

\subsection{Борьба за ренту (Таллок, 1980)}

Модель состязательного распределения ресурсов применяется для описания спорта, гонки вооружений, НИОКР, судебных тяжб и лоббирования. Две фирмы соревнуются за право получения приза (например, постройка магазина на центральной площади города) размером $R$. Успех не гарантирован, однако его вероятность возрастает с ростом невозвратных инвестиций (издержек на лоббирование) $r_i \in [0, \infty)$.

Вероятность успеха фирмы $i \in \{1, 2\}$ задается функцией:
\[
P_i = \frac{r_i^\gamma}{r_1^\gamma + r_2^\gamma}
\]
где $\gamma \ge 0$ — параметр, отражающий эффективность лоббирования. Предельными случаями являются $\gamma = 0$ (вероятность не зависит от усилий) и $\gamma = \infty$ (полная детерминированность: выигрывает тот, кто вложил больше).

Ожидаемый выигрыш фирмы $i$ составляет:
\[
u_i(r_1, r_2) = R \frac{r_i^\gamma}{r_1^\gamma + r_2^\gamma} - r_i
\]

Если профиль $(r_1, r_2)$ является внутренним равновесием ($r_i > 0$), должны выполняться условия первого порядка (FOC): $\frac{\partial u_i}{\partial r_i} = 0$. Угловые решения характеризуются условием $\frac{\partial u_i}{\partial r_i} \le 0$ при $r_i = 0$.

Вычислим частные производные:
\begin{align}
\frac{\partial u_1}{\partial r_1} &= R \frac{\gamma r_1^{\gamma-1} r_2^\gamma}{(r_1^\gamma + r_2^\gamma)^2} - 1 \label{eq:tullock_foc1} \\
\frac{\partial u_2}{\partial r_2} &= R \frac{\gamma r_2^{\gamma-1} r_1^\gamma}{(r_1^\gamma + r_2^\gamma)^2} - 1 \label{eq:tullock_foc2}
\end{align}

Если $r_2 = 0$, то $\frac{\partial u_1}{\partial r_1} = -1 \le 0$, следовательно, фирма 1 выберет $r_1 = 0$. Однако профиль $(0, 0)$ не является равновесием, так как бесконечно малое положительное отклонение $r_1' > 0$ принесет фирме 1 почти весь приз $R$. Следовательно, условия первого порядка обязаны выполняться со знаком равенства для строго положительных усилий.

Из симметрии уравнений \eqref{eq:tullock_foc1} и \eqref{eq:tullock_foc2} следует равновесный профиль:
\[
r_1 = r_2 = R \frac{\gamma}{4}
\]
(При нормировке приза $R = 1$ равновесные стратегии принимают вид $r_1 = r_2 = \frac{\gamma}{4}$). 

Существование равновесия зависит от параметра $\gamma$:
\begin{itemize}
    \item При $\gamma \le 1$: функции полезности глобально квазивогнуты, найденное равновесие существует.
    \item При $\gamma \in (1, 2]$: функции полезности не являются квазивогнутыми, однако локальный максимум остается глобальным, равновесие существует.
    \item При $\gamma > 2$: условия второго порядка нарушаются, функции полезности сильно выпуклы на начальном участке, равновесия в чистых стратегиях нет (фирмам выгодно отклониться в 0).
\end{itemize}


\subsection{Модель линейного города (H. Hotelling, 1929)}

Рассматривается город в виде отрезка длиной $l$. Две фирмы продают однородный товар. Фирма $A$ расположена на расстоянии $a$ от левого края, фирма $B$ — на расстоянии $b$ от правого края. Предполагается, что $a + b \le l$. 

Потребители равномерно распределены вдоль города с общей плотностью, нормированной к 1. Каждый потребитель покупает ровно одну единицу товара. Транспортные издержки линейны: перемещение на $t$ метров обходится потребителю в $ct$ денежных единиц.

Определим положение безразличного потребителя, находящегося между фирмами на расстоянии $x$ от фирмы $A$ и на расстоянии $y$ от фирмы $B$. Расстояния связаны соотношением $y = l - a - b - x$. Условие безразличия (равенство полных затрат на покупку у фирмы $A$ и фирмы $B$) имеет вид:
\[
p_1 + cx = p_2 + c y \implies p_1 + cx = p_2 + c(l - a - b - x)
\]
Выражая $x$ и $y$, получаем величину спроса для каждой фирмы ($q_1 = a + x$ и $q_2 = b + y$):
\begin{align*}
x &= \frac{1}{2} \left( l - a - b + \frac{p_2 - p_1}{c} \right) \\
y &= \frac{1}{2} \left( l - a - b + \frac{p_1 - p_2}{c} \right)
\end{align*}

Функции полезности (прибыли) фирм без учета издержек производства задаются как $u_1 = p_1 q_1$ и $u_2 = p_2 q_2$:
\begin{align*}
u_1(p_1, p_2) &= \frac{1}{2} (l + a - b) p_1 - \frac{p_1^2}{2c} + \frac{p_1 p_2}{2c} \\
u_2(p_1, p_2) &= \frac{1}{2} (l - a + b) p_2 - \frac{p_2^2}{2c} + \frac{p_1 p_2}{2c}
\end{align*}

Максимизация прибыли ($\frac{\partial u_1}{\partial p_1} = 0$ и $\frac{\partial u_2}{\partial p_2} = 0$) дает систему условий первого порядка:
\[
\begin{cases} 
\frac{1}{2}(l + a - b) - \frac{p_1}{c} + \frac{p_2}{2c} = 0 \\
\frac{1}{2}(l - a + b) - \frac{p_2}{c} + \frac{p_1}{2c} = 0 
\end{cases}
\]

Разрешая систему относительно цен, находим равновесие Нэша в ценах:
\begin{align*}
p_1 &= c \left( l + \frac{a - b}{3} \right), & p_2 &= c \left( l - \frac{a - b}{3} \right)
\end{align*}
Равновесные объемы продаж:
\begin{align*}
q_1 &= \frac{1}{2} \left( l + \frac{a - b}{3} \right), & q_2 &= \frac{1}{2} \left( l - \frac{a - b}{3} \right)
\end{align*}
Равновесные прибыли:
\begin{align*}
u_1 &= \frac{c}{2} \left( l + \frac{a - b}{3} \right)^2, & u_2 &= \frac{c}{2} \left( l - \frac{a - b}{3} \right)^2
\end{align*}

Анализ функций прибыли показывает, что у обеих фирм существует стимул увеличивать свои координаты $a$ и $b$, двигаясь навстречу друг другу вплоть до положения $a + b = l$. Этот результат известен как \textbf{принцип минимальной дифференциации}: фирмы стремятся располагаться максимально близко друг к другу (в центре города при условии симметрии).

\subsection{Проблемы классической модели Хотеллинга}

Классический анализ Хотеллинга базировался на предпосылке о том, что безразличный потребитель всегда находится строго между фирмами ($x \ge 0$ и $y \ge 0$). В работе \textit{D’Aspremont, Gabszewicz, Thisse (1979)} было показано, что этот вывод ошибочен.

Если фирмы располагаются в одной точке (или очень близко друг к другу), возникает ситуация совершенной конкуренции по Бертрану. Одна фирма может незначительно снизить цену, тем самым «подрезав» конкурента и захватив весь рынок. В этом случае равновесные цены и прибыли падают до нуля, что означает наличие стимула к удалению друг от друга.

С учетом возможности захвата всего рынка, полные функции полезности имеют кусочно-заданный вид:
\[
u_1(p_1, p_2) = 
\begin{cases} 
\frac{1}{2} (l + a - b) p_1 - \frac{p_1^2}{2c} + \frac{p_1 p_2}{2c}, & |p_1 - p_2| \le c(l - a - b) \\
l p_1, & p_1 < p_2 - c(l - a - b) \\
0, & p_1 > p_2 + c(l - a - b)
\end{cases}
\]
\[
u_2(p_1, p_2) = 
\begin{cases} 
\frac{1}{2} (l - a + b) p_2 - \frac{p_2^2}{2c} + \frac{p_1 p_2}{2c}, & |p_1 - p_2| \le c(l - a - b) \\
l p_2, & p_2 < p_1 - c(l - a - b) \\
0, & p_2 > p_1 + c(l - a - b)
\end{cases}
\]

Функция прибыли обладает двумя локальными максимумами. Анализ показывает, что равновесие, найденное Хотеллингом, существует только при выполнении строгих условий на локацию:
\[
\begin{cases}
\left( l + \frac{a - b}{3} \right)^2 \ge \frac{4}{3} l (a + 2b) \\
\left( l + \frac{b - a}{3} \right)^2 \ge \frac{4}{3} l (b + 2a)
\end{cases}
\]
При требовании симметричности локаций ($a = b$) равновесие в ценах существует лишь в том случае, если фирмы располагаются достаточно далеко друг от друга: $a = b \le \frac{l}{4}$. Если добавить полную свободу локаций в качестве первой стадии игры, то из-за тенденции к агломерации в центре ценовое равновесие на второй стадии разрушается, следовательно, чистого равновесия в полной игре нет, а принцип минимальной дифференциации не работает.

\subsection{Модификация с квадратичными издержками}

Для восстановления существования равновесия вводится предположение о выпуклых (квадратичных) транспортных издержках: перемещение на $t$ метров стоит $ct^2$ рублей.

Структура функции полезности остается схожей, однако условия на точки разрыва меняются таким образом, что «центральный пик» функции прибыли всегда оказывается выше боковых для любых локаций. Равновесие в ценах существует всегда и задается формулами:
\begin{align*}
p_1 &= c (l - a - b) \left( l + \frac{a - b}{3} \right) \\
p_2 &= c (l - a - b) \left( l + \frac{b - a}{3} \right)
\end{align*}

Подстановка этих цен в функции прибыли показывает, что каждая равновесная полезность строго убывает с ростом собственной локации ($a$ или $b$). Таким образом, у фирм возникают стимулы максимально удалиться друг от друга, что приводит к \textbf{принципу максимальной дифференциации}.

Этот результат демонстрирует, что в теории пространственной конкуренции не существует единой глобальной закономерности. Выбор конкретной спецификации модели (в частности, функции транспортных издержек) имеет принципиальное значение. Дальнейшие модификации включают:
\begin{itemize}
    \item Модель кругового города (Salop, 1979) для анализа конкуренции с соседями и свободного входа множества фирм.
    \item Модели с многомерным пространством и различными метриками (например, метрика Манхэттена).
    \item Модель «спицевидного города» (Chen, Riordan, 2007) для исследования влияния входа на рынок.
\end{itemize}
